\chapter{Dificuldades e decisões técnicas}
\label{cap:dificuldades}

\section{Correspondência semântica com Drools}
\label{sec:dif-drools}

A dificuldade principal foi reproduzir em Prolog comportamentos que o Drools
oferece diretamente, como memória com estado, eventos com expiração, janelas
temporais, \emph{salience} e reatividade a atualizações. A solução separou o
estado persistente da memória do ciclo, codificou as janelas por
\emph{timestamps} e criou um motor com agenda explícita.

\section{Ordem dos eventos}
\label{sec:dif-eventos}

Eventos relacionados podem chegar em ordens diferentes. O
\cod{driving_state.pl} processa as distâncias frontais antes das restantes
medições para que a velocidade relativa e o TTC atualizem o alvo correto. Este
detalhe evita decisões dependentes da ordem incidental do \emph{array} JSON.

\section{Explicação de condições auxiliares}
\label{sec:dif-auxiliares}

Um \cod{test/1} comum pode ocultar os factos consultados por um predicado
auxiliar. A forma \cod{prove/1} resolve este problema: prova o objetivo por
encadeamento inverso, recolhe os identificadores dos factos lidos e permite
construir uma cadeia explicativa. A mesma decomposição torna possível explicar
por que uma mudança de faixa não é segura.

\section{Conflitos entre velocidade e faixa}
\label{sec:dif-canais}

Uma única ordenação global faria uma proposta de mudança de faixa competir
diretamente com uma travagem, embora controlem dimensões diferentes. O
\cod{DecisionContext} divide as propostas em canal longitudinal e lateral e
aplica depois um veto de segurança. Esta modelação permite travar enquanto se
inicia uma manobra quando isso é seguro, mas bloqueia a manobra perante uma
paragem obrigatória mais importante.
